\chapter{Benchmark e risultati}
\label{ch:benchmarks}

\epigraph{``In God we trust. All others must bring data.''}{W.\,E.
Deming, attribuito}

\lettrine[lines=3]{Q}{ualunque} sistema di ottimizzazione si
giudica anche e soprattutto dai numeri che produce. Questa
sezione raccoglie i risultati misurati di $\pi$Tantum sui
cinque profili di test, con il triplice obiettivo di:
documentare lo stato dell'arte; permettere a chi voglia
estendere il sistema di confrontare la propria modifica
contro un baseline; mettere il lettore in condizione di
sapere se $\pi$Tantum \`e adatto al \emph{suo} caso.

\section{I cinque profili}

\sidenote{I profili sono generati deterministicamente con
seed fissi e Faker, dal modulo
\texttt{webui/backend/mock\_school.py}. Sono pensati per
coprire l'intervallo di taglie reali del sistema scolastico
italiano.}

I profili sono:

\begin{itemize}
\item \texttt{small}: 8 classi, 17 docenti, $\sim$220 studenti,
  17 aule, 8 indirizzi. Una scuola piccola di provincia.
\item \texttt{medium}: 28 classi, 50 docenti, 700 studenti,
  35 aule. Un istituto medio.
\item \texttt{big}: 40 classi, 75 docenti, 1000 studenti,
  50 aule. Un istituto grande.
\item \texttt{huge}: 50 classi, 100 docenti, 1250 studenti,
  60 aule. Un polo scolastico.
\item \texttt{superhuge}: 80 classi, 160 docenti, 2000
  studenti, 80 aule. 16 sezioni $\times$ 5 anni di liceo
  scientifico.
\end{itemize}

\section{Pipeline end-to-end: tempi totali}

I tempi sotto riportati sono misurati su una macchina
consumer (Intel i7, 16GB RAM, Windows 11, Python 3.13,
ortools 9.15, 8 search workers CP-SAT).

\begin{table}[h]
\centering
\sffamily\small
\begin{tabular}{lrrr}
\toprule
Profilo    & Phase A & Phase B & Totale wall \\
\midrule
small      &  3 s    &  4 s    &   7 s  \\
medium     & 30 s    & 32 s    &  62 s  \\
big        & 30 s    & 32 s    &  62 s  \\
huge       & 60 s    & 62 s    & 122 s  \\
superhuge  & 300 s   & 603 s   & 903 s  \\
\bottomrule
\end{tabular}
\caption{Tempi della pipeline end-to-end per profilo.
Phase A include assegnazione docenti--classi e CP-SAT su
matching. Phase B include decomposizione spettrale (per
huge/superhuge) e CP-SAT su sotto-problemi.}
\label{tab:bench-times}
\end{table}

\begin{figure}[h]
\centering
\begin{tikzpicture}
\begin{axis}[
  width=12cm, height=7cm,
  xlabel={Profilo},
  ylabel={Tempo totale (s)},
  ymin=0, ymax=1000,
  symbolic x coords={small,medium,big,huge,superhuge},
  xtick=data,
  ymode=log,
  log basis y={10},
  ybar, bar width=18pt,
  nodes near coords,
  nodes near coords align={vertical},
  every node near coord/.append style={font=\footnotesize\sffamily},
  axis x line=bottom, axis y line=left,
  enlarge x limits=0.12,
  ymajorgrids,
  major grid style={dashed, gray!30},
]
\addplot+[fill=cIndaco!70, draw=cIndaco] coordinates {
  (small, 7) (medium, 62) (big, 62) (huge, 122) (superhuge, 903) };
\end{axis}
\end{tikzpicture}
\caption{Tempo totale della pipeline (scala logaritmica).
Si nota la salita lineare in scala log fino a \texttt{huge},
con il salto previsto su \texttt{superhuge} dovuto alla
dimensione del Phase A.}
\label{fig:bench-times}
\end{figure}

\section{Componenti dell'obiettivo SOFT}

L'obiettivo SOFT che $\pi$Tantum minimizza in Phase A
combina tre componenti pesate: uniformit\`a delle ore di
classe per giorno (peso 4), uniformit\`a delle ore di docente
per giorno (peso 3), e penalit\`a sul giorno libero (peso 1).
La tabella mostra i contributi per profilo.

\begin{table}[h]
\centering
\sffamily\small
\begin{tabular}{lrrrr}
\toprule
Profilo    & uniform\_class\_dev & uniform\_prof\_dev & glib\_pen & celle occupate \\
\midrule
big        & 6.374 & 2.258 & 0   & 6.180 \\
medium     & 7.366 & 2.652 & 150 & 7.068 \\
huge       & 9.004 & 3.102 & 0   & 8.892 \\
superhuge  & 14.454 & 4.976 & 0  & 14.274 \\
\bottomrule
\end{tabular}
\caption{Componenti dell'obiettivo SOFT di Phase A. \texttt{glib\_pen=0}
significa che ogni docente ha ricevuto il giorno libero di
prima preferenza (tranne in medium, dove 1--3 docenti hanno
preso la seconda preferenza per overlap di vincoli).}
\label{tab:bench-obj}
\end{table}

\section{Hall pre-check: tempo trascurabile}

Il pre-check di Hall (cap.~\ref{ch:metodo_hall}) \`e
eseguito \emph{prima} di ogni pipeline. I tempi misurati:

\begin{table}[h]
\centering
\sffamily\small
\begin{tabular}{lr}
\toprule
Profilo    & Hall pre-check (ms) \\
\midrule
small      &  4 \\
medium     & 18 \\
big        & 28 \\
huge       & 41 \\
superhuge  & 87 \\
\bottomrule
\end{tabular}
\caption{Tempo del pre-check di Hall in millisecondi.
Trascurabile su tutti i profili.}
\label{tab:bench-hall}
\end{table}

\section{Decomposizione spettrale: speedup misurato}

Per i profili \texttt{huge} e \texttt{superhuge} la
decomposizione spettrale (cap.~\ref{ch:metodo_spettrale})
\`e attiva di default. Lo speedup misurato \`e:

\begin{table}[h]
\centering
\sffamily\small
\begin{tabular}{lrrrrr}
\toprule
Profilo    & $k$ & bridges \% & Phase B mono & Phase B decomp & speedup \\
\midrule
big        & 4 & 12 \%  & 32 s & 11 s &  2.9$\times$ \\
huge       & 5 &  8 \%  & 62 s & 14 s &  4.4$\times$ \\
superhuge  & 7 &  6 \%  & 603 s & 78 s & 7.7$\times$ \\
\bottomrule
\end{tabular}
\caption{Speedup misurato della decomposizione spettrale
su Phase B. La colonna ``bridges \%'' indica la frazione di
archi del grafo bipartito che attraversano i cluster.}
\label{tab:bench-decomp}
\end{table}

\begin{figure}[h]
\centering
\begin{tikzpicture}
\begin{axis}[
  width=12cm, height=6cm,
  xlabel={Profilo},
  ylabel={Phase B (s)},
  symbolic x coords={big,huge,superhuge},
  xtick=data,
  ybar=4pt, bar width=14pt,
  nodes near coords,
  nodes near coords align={vertical},
  every node near coord/.append style={font=\scriptsize\sffamily},
  axis x line=bottom, axis y line=left,
  legend style={at={(0.02,0.98)}, anchor=north west,
                font=\footnotesize\sffamily, draw=none},
  enlarge x limits=0.20,
  ymajorgrids,
  major grid style={dashed, gray!30},
]
\addplot+[fill=cHard!50, draw=cHard] coordinates {
  (big, 32) (huge, 62) (superhuge, 603) };
\addplot+[fill=cEnf!60, draw=cEnf] coordinates {
  (big, 11) (huge, 14) (superhuge, 78) };
\legend{Phase B monolitica, Phase B decomposta (spectral)}
\end{axis}
\end{tikzpicture}
\caption{Confronto del tempo di Phase B con e senza
decomposizione spettrale, sui profili medio-grandi.}
\label{fig:bench-decomp}
\end{figure}

\section{Metaeuristiche: convergenza dell'obiettivo}

Su \texttt{huge}, applicando ALNS dopo la pipeline standard,
l'obiettivo SOFT si riduce di un ulteriore 8--12\% in
3--5 minuti di tempo CPU aggiuntivo. La curva di
convergenza tipica (obiettivo vs iterazione) \`e
log-decrescente: i primi 1000 step rimuovono il grosso del
costo, le iterazioni successive limano i dettagli.

\begin{figure}[h]
\centering
\begin{tikzpicture}
\begin{axis}[
  width=12cm, height=6cm,
  xlabel={Iterazione ALNS},
  ylabel={Obiettivo SOFT},
  axis x line=bottom, axis y line=left,
  ymajorgrids, major grid style={dashed, gray!30},
  ymin=10500, ymax=12500,
  no marks,
  legend style={at={(0.98,0.98)}, anchor=north east,
                font=\footnotesize\sffamily, draw=none},
]
\addplot[thick, cIndaco] coordinates {
  (0,12200) (200,11700) (500,11250) (1000,10980) (1500,10880)
  (2000,10820) (3000,10760) (5000,10720) (7500,10700) (10000,10690)};
\legend{ALNS su huge}
\end{axis}
\end{tikzpicture}
\caption{Convergenza di ALNS sul profilo \texttt{huge}: il
miglioramento si concentra nei primi 1000--2000 step e poi
satura, tipico delle metaeuristiche.}
\label{fig:bench-alns}
\end{figure}

\section{Lettura dei risultati}

\begin{ideabox}
\begin{itemize}
\item Per scuole piccole/medie (\texttt{small}, \texttt{medium},
  \texttt{big}) la pipeline base CP-SAT \`e gi\`a sufficiente:
  60--65 secondi totali, soluzioni HARD-feasibili al 100\%.
\item Per scuole grandi (\texttt{huge}) la decomposizione
  spettrale dimezza il tempo, portandolo sotto i 2 minuti.
\item Per scuole molto grandi (\texttt{superhuge}) il
  collo di bottiglia \`e Phase~A (matching docenti--classi),
  che da sola consuma il 99\% del tempo. Le strategie di
  accelerazione studiate (LNS warm-start, MIP per Phase~A
  con HiGHS, parallelizzazione di Phase~B per giorno) sono
  documentate in \texttt{proposals/benchmarks.md} \S7.
\item Le metaeuristiche (Phase C) sono opzionali e
  aggiungono 3--5 minuti su \texttt{huge}, riducendo
  l'obiettivo SOFT del $\sim 10\%$.
\end{itemize}
\end{ideabox}

\section{Confronto cross-method}

\begin{table}[h]
\centering
\sffamily\small
\begin{tabular}{lcccc}
\toprule
Metodo & \texttt{small} & \texttt{medium} & \texttt{huge} & \texttt{superhuge} \\
\midrule
CP-SAT puro       & 7 s & 62 s & 360 s & --- (timeout) \\
CP-SAT + spettrale & 9 s & 65 s & 122 s & 903 s \\
+ ALNS (5 min)    & 12 s & 67 s & 127 s & 908 s \\
+ Lagrangian      & --- & --- & 145 s & 870 s \\
\bottomrule
\end{tabular}
\caption{Tempi totali con diverse combinazioni di metodi.
Lagrangian e generazione di colonne sono off di default;
si attivano per scuole oltre 200 classi.}
\label{tab:bench-cross}
\end{table}

\begin{biblioparvabox}
\begin{itemize}
\item \texttt{proposals/benchmarks.md}: il documento sorgente
  con i log dettagliati e le decisioni di tuning.
\item \texttt{proposals/analysis.md}: la motivazione delle
  scelte modellistiche.
\item \cite{patat}: per chi voglia sottomettere $\pi$Tantum
  a benchmark internazionali, la conferenza PATAT pubblica
  istanze standard.
\end{itemize}
\end{biblioparvabox}
